\section{Introduction}

In electric machine control, minimizing copper losses while maintaining a
desired torque is a fundamental objective. The Maximum Torque Per Ampere
(MTPA) strategy aims to find stator currents \(\mathbf{x}\) that produce the
required torque with minimum copper losses, represented by the quadratic form
\(\mathbf{x}^\top Q\mathbf{x}\), where \(Q\) is typically a diagonal matrix
containing the stator resistances.

Traditional quadratic programming approaches solve such problems globally
but may not easily handle constraints defined by the convex hull of feasible
current sets, especially when the constraint domain is a simplex defined by a
finite set of vertices. In this paper, we propose a recursive quadratic
programming approach that exploits the simplex structure of the feasible set
and recursively reduces dimensionality to find the global optimum within the
simplex.

\section{Problem Formulation}

Consider the quadratic minimization problem
\begin{equation}
  \min_{\mathbf{x}} \quad \mathbf{x}^\top Q\mathbf{x}
  \quad \text{s.t.} \quad
  \mathbf{x}\in\operatorname{conv}(\mathbf{v}_1,\mathbf{v}_2,\ldots,\mathbf{v}_m),
\end{equation}
where \(Q\in\mathbb{R}^{n\times n}\) is symmetric positive semidefinite and
the feasible set is the convex hull of \(m\) vertices
\(\mathbf{v}_i\in\mathbb{R}^n\). The vertices represent extremal current
vectors satisfying operational constraints.

Because the feasible region is a simplex, any point \(\mathbf{x}\) in the
convex hull can be expressed as
\[
  \mathbf{x}=\mathbf{v}_1+B\boldsymbol{\lambda},
\]
where \(B=[\mathbf{v}_2-\mathbf{v}_1,\ldots,\mathbf{v}_m-\mathbf{v}_1]\),
\(\boldsymbol{\lambda}\geq0\), and \(\sum\lambda_i\leq1\).

\section{Recursive Solution Algorithm}

The recursive approach proceeds by attempting to solve the unconstrained
quadratic minimization in barycentric coordinates
\(\boldsymbol{\lambda}\) on the simplex, checking feasibility, and, if the
candidate lies outside, recursively solving on lower-dimensional faces. This
ensures that the global optimum lies within or on the boundary of the
simplex.

\begin{lstlisting}[style=matlab,caption={Recursive quadratic minimization within a simplex},label={lst:recursive_qp}]
function x_optimal = recursive_QP_within_simplex(vertices, Q)
% Recursive solver for min x'Qx subject to x in conv(vertices)

if size(vertices, 2) >= 2
    B = vertices(:, 2:end) - vertices(:, 1);
    QBB = B' * Q * B;
    QBv = B' * Q * vertices(:, 1);

    if rcond(QBB) < eps
        is_inside = false;
    else
        lambda = -QBB \ QBv;
        is_inside = all(lambda >= -eps & lambda <= 1+eps) ...
            && (sum(lambda) <= 1+eps);
    end

    if is_inside
        x_optimal = vertices(:, 1) + B * lambda;
    else
        subset_indices = nchoosek(1:size(vertices, 2), ...
            size(vertices, 2) - 1);
        x_candidates = NaN(size(vertices, 1), size(subset_indices, 1));
        f_values = NaN(1, size(subset_indices, 1));

        for idx = 1:size(subset_indices, 1)
            subset = vertices(:, subset_indices(idx, :));
            x_sub = recursive_QP_within_simplex(subset, Q);
            x_candidates(:, idx) = x_sub;
            f_values(idx) = x_sub' * Q * x_sub;
        end
        [~, best_idx] = min(f_values);
        x_optimal = x_candidates(:, best_idx);
    end
else
    x_optimal = vertices;
end
end
\end{lstlisting}

\section{Application to Maximum Torque per Current}

In the MTPA problem, the vector \(\mathbf{x}\) contains the stator currents
in each phase or the $dq$-axis components, and
\(Q=\operatorname{diag}(R_1,R_2,\ldots,R_n)\) represents the resistances
responsible for copper losses. Minimizing \(\mathbf{x}^\top Q\mathbf{x}\)
corresponds to minimizing these losses for a given torque, where the feasible
set is constrained by current limits and torque requirements that can be
represented by the convex hull of vertices.

The recursive quadratic programming method guarantees that the resulting
solution respects the convex constraints while minimizing copper losses
efficiently, even when the number of vertices is large or the problem is
defined in higher dimensions.

\section{Conclusion}

We have presented a recursive quadratic programming solver tailored for
optimization problems constrained to simplices, with application to MTPA
current optimization in electric machines. The method recursively reduces
dimensionality to handle boundary conditions and guarantees the global
minimum within the convex hull. Future work may focus on extending this
method to more complex nonlinear constraints and real-time implementations.
